\chapter{体会}
\label{chap:reflection}

\section{项目层面的工程认识}
本节总结可由源码和测试印证的工程认识，不冒称为某位成员的个人经历。个人完成了什么、何时遇到问题、怎样协作等事实，需要各成员在后续小节独立补充。

\subsection{业务规则需要明确的状态边界}
同一页面同时包含真人与虚拟对象，容易把演示操作误描述为所有用户共有能力。将纸条、聊天、心笺及授权按对象类型分析，才能正确确定谁能够执行操作、哪些状态允许转换以及失败时是否消耗额度。服务端白名单、数据库约束和测试中的非法动作前后状态比较说明：可靠交互不能只靠按钮置灰，还需要在写入位置执行规则。

\subsection{异步结果不等于当前有效结果}
模型生成与用户继续输入、切换风格或修改图片可能同时发生。消息编号、当前风格、参考图摘要和目标图片版本提供了判断结果是否仍适用的条件。文本分析可以丢弃过期响应；图片生成还必须经过私有预览和确认保存。两类处理共同体现了一个原则：上游请求成功并不能跳过业务一致性检查。

\subsection{降级和恢复应成为设计的一部分}
真人发言不依赖模型代回，普通动作与模型并发额度分离；虚拟回复失败保留任务，重试不重复写消息；生成失败不覆盖已有图。这些实现将失败纳入正常控制流，而非只提供一个通用报错。数据库事务和重启状态处理使恢复条件可解释，但仍不能推导出已经具有生产级灾难恢复能力。

\subsection{证据的适用范围需要随版本维护}
早期文档的“没有账号、没有真人通信”不能无日期地用于当前版本；历史测试环境和性能也不能改标签后当作本次结果。源码行数可以说明项目规模，却不能说明成员贡献。报告将当前源码、模拟测试、历史测量、截图和未取得的用户反馈分开，有助于发现描述漂移，并避免把课程演示证据过度推广。

\section{成员个人体会（待补充）}
当前未取得成员名单和实际经历，不能据模板四行槽位认定有四位开发者。实际人数确定后，每位成员应分别填写以下内容，并与第\ref{chap:problem}章的实际分工及提交或任务记录一致。
\begin{enumerate}
  \item 姓名、学号与承担任务：待填写。说明负责的具体模块、交付物及合作边界。
  \item 遇到的实际困难：待填写。描述触发条件、错误现象、失败尝试和当时限制，不使用泛泛的“学会合作”。
  \item 分析与解决过程：待填写。列出定位依据、方案取舍、本人实施内容和协作对象。
  \item 验证与结果：待填写。提供测试、提交、截图或复测证据，说明是否彻底解决及残余问题。
  \item 反思与后续改进：待填写。说明如果重做，会怎样调整需求、接口、验证或时间安排。
\end{enumerate}

本节不是虚构的组员自述，也不把项目层面的技术总结平均归属于成员。填写时应保留具体问题与解决方法，避免抄录相同体会或用代码行数代替贡献说明。
